\باب{ریاضیاتی ترغیب}\شناخت{ضمیمہ_الف}

\اصطلاح{ریاضیاتی  ترغیب کا اصول}\فرہنگ{ریاضیاتی ترغیب}\حاشیہب{mathematical induction}\فرہنگ{mathematical induction} استعمال کر کے   ثابت کیا جا سکتا ہے کہ  کئی کلیات ہر مثبت عدد صحیح کے لئے مطمئن ہوں گے۔ ایسا ایک کلیہ درج ذیل ہے۔
\[1+2+\cdots+n=\frac{n(n+1)}{2}\]
جو ثبوت اس مسلمہ  کو استعمال کرتا ہے وہ  \ترچھا{ترغیبی ثبوت }کہلاتا ہے۔

ترغیب سے کلیے کا ثبوت درج ذیل اقدام سے کیا جاتا ہے۔
\begin{enumerate}[1.]
\item
تصدیق کریں کہ کلیہ \عددی{n=1} کے لئے درست ہے۔
\item
ثابت کریں کہ  اگر کلیہ کسی مثبت عدد صحیح  \عددی{n=k} کے لئے درست ہو تب یہ \عددی{n=k+1} کے لئے بھی درست ہو گا۔
\end{enumerate}

مسلمہ ترغیب کہتا ہے ان دو اقدام  کے مکمل ہونے پر کلیہ تمام مثبت عدد صحیح \عددی{n}  کے لئے درست ہو گا۔قدم  اول تصدیق کرتا ہے کہ کلیہ \عددی{n=1} کے لئے درست ہے۔  قدم    دوم کے تحت اگر کلیہ  \عددی{n=2} کے لئے درست ہو تب  کلیہ  \عددی{n=3} کے لئے بھی درست ہو گا، اور اسی طرح قدم دوم  کے تحت اگر  کلیہ  \عددی{n=3} کے لئے درست ہو تب یہ \عددی{n=4} کے لئے بھی درست ہو گا، وغیرہ۔

\ابتدا{مثال}
درج ذیل کلیہ ہر مثبت عدد صحیح \عددی{n} کے لئے درست ثابت کریں۔
\[1+2+\cdots+n=\frac{n(n+1)}{2}\]

حل:\quad
ہم  مذکورہ بالا دو اقدام پر چلتے ہوئے ثبوت پیش کرتے  ہیں۔
\begin{enumerate}[1.]
\item
ہم کلیہ میں دونوں اطراف \عددی{n=1} پُر کر کے قدم اول  کی تصدیق  کرتے ہیں۔
\[1=\frac{1(1+1)}{2}\]
\item
قدم دوم میں ہم فرض کرتے ہیں کلیہ  مثبت عدد صحیح \عددی{n=k} کے لئے درست ہے، یعنی ہم مان لیتے ہیں کہ درج ذیل درست ہے۔
\[1+2+\cdots+k=\frac{k(k+1)}{2}\]
اب اگلا  مثبت عدد صحیح \عددی{k+1} ہو گا لہٰذا  ہم کلیے کا صرف بایاں ہاتھ لے کر اس میں \عددی{k+1} جمع کر کے حل کرتے ہیں۔
\begin{align*}
\underbrace{1+2+\cdots+k}_{\frac{k(k+1)}{2}}+(k+1)&=\frac{k(k+1)}{1}+(k+1)=\frac{k^2+k+2k+2}{2}\\
&=\frac{(k+1)(k+2)}{2}=\frac{(k+1)(k+1)+1}{2}
\end{align*}
یوں دونوں اطراف \عددی{k+1} کی جگہ \عددی{n} لکھ کر درج ذیل حاصل ہو گا۔
\[1+2+\cdots+n=\frac{n(n+1)}{2}\]
یعنی کلیے کو \عددی{n=k} کے لئے درست مان کر ہم نے ثابت کیا کہ یہی کلیہ \عددی{n=k+1} کے لئے بھی درست ہے۔

یوں ریاضیاتی ترغیب کا اصول  اس بات کی ضمانت دیتا ہے کہ تمام مثبت عدد صحیح \عددی{n} کے لئے دیا گیا کلیہ درست ہے۔
\end{enumerate}
\انتہا{مثال}

حصہ \حوالہء{5.2}  کی مثال \حوالہء{4} میں  اس کلیے کا دوسرا ثبوت پیش کیا گیا  تھا۔ تاہم، ریاضیاتی ترغیب سے  ثبوت کا حصول زیادہ عمومی طریقہ  ہے۔ اس سے  ابتدائی \عددی{n} عدد صحیح   کے مربعوں  اور مکعب کے مجموعوں  کے کلیات ثابت کیے جا سکتے ہیں (مشق \حوالہء{9} اور \حوالہء{10})۔
دوسری   مثال پیش کرتے ہیں۔

\ابتدا{مثال}
ریاضیاتی ترغیب استعمال کر کے ثابت کریں کہ تمام مثبت عدد صحیح کے لئے درج ذیل کلیہ درست ہو گا۔
\[\frac{1}{2^1}+\frac{1}{2^2}+\cdots+\frac{1}{2^n}=1-\frac{1}{2^n}\]

حل:\quad
ہم ریاضیاتی ترغیب کی دو اقدام مکمل کر کے ثبوت پیش کرتے ہیں۔
\begin{enumerate}[1.]
\item
ہم کلیہ \عددی{n=1} کے لئے جانچتے ہیں۔
\[\frac{1}{2^1}=1-\frac{1}{2^1}\]
\item
اگر \عددی{n=k} کے لئے درج ذیل درست تسلیم کیا جائے
\[\frac{1}{2^1}+\frac{1}{2^2}+\cdots+\frac{1}{2^k}=1-\frac{1}{2^k}\]
تب بائیں ہاتھ اگلا جزو جمع کر کے حل کرتے ہیں۔
\begin{align*}
\frac{1}{2^1}+\frac{1}{2^2}+\cdots+\frac{1}{2^k}+\frac{1}{2^{k+1}}&=1-\frac{1}{2^k}+\frac{1}{2^{k+1}}=1-\frac{1\cdot 2}{2^k\cdot 2}+\frac{1}{2^{k+1}}\\
&=1-\frac{2}{2^{k+1}}+\frac{1}{2^{k+1}}=1-\frac{1}{2^{k+1}}
\end{align*}
یوں  اگر دیا  گیا کلیہ \عددی{n=k}  کے لئے درست ہو تب یہ  \عددی{n=k+1} کے لئے بھی درست ہو گا۔

ان اقدام کی تصدیق کے بعد اصول ریاضیاتی ترغیب ضمانت دیتا ہے کہ کلیہ ہر مثبت عدد صحیح کے لئے درست ہو گا۔
\end{enumerate}
\انتہا{مثال}

\جزوحصہء{دیگر ابتدائی عدد صحیح}
بعض  ترغیبی  دلائل میں \عددی{n=1} سے آغاز کرنے کے بجائے کسی  اور عدد صحیح  سے  ابتدا کی جاتی ہے۔ اس طرح کے دلائل کے قدم درج ذیل ہیں۔ 

\begin{enumerate}[1.]
\item
تصدیق کریں کہ کلیہ \عددی{n=n_1} (متعلقہ پہلا عدد صحیح)  کے لئے درست ہے۔
\item
ثابت کریں کہ اگر کلیہ کسی بھی عدد صحیح  \عددی{n=k\ge n_1}  کے لئے درست ہو تب یہ \عددی{n=(k+1)} کے لئے بھی درست ہو گا۔
\end{enumerate}
ان اقدام کے مکمل ہونے کے بعد اصول ریاضیاتی ترغیب ضمانت دیتا ہے کہ کلیہ تمام \عددی{n\ge n_1} کے لئے درست ہو گا۔

\ابتدا{مثال}
دکھائیں کہ اگر \عددی{n} کافی بڑا ہو تب \عددی{n!>3^n} ہو گا۔

حل:\quad
یہاں کتنا بڑا عدد صحیح کافی بڑا تصور کیا جا سکتا ہے؟ یہ جانتے ہیں۔
\[
\begin{array}{l|rrrrrrr}
n&1&2&3&4&5&6&7\\
\midrule
n!&1&2&6&24&120&720&5040\\
\midrule
3^n&3&9&27&81&243&729&2187
\end{array}
\]
ان نتائج کو دیکھ کر  \عددی{n!>3^n}  کے لئے \عددی{n\ge 7} ہونا   کافی ہے۔ اس کی مکمل  تصدیق کے لئے ہم ریاضیاتی ترغیب بروئے کار لاتے ہیں۔ ہم قدم اول میں \عددی{n_1=7} منتخب کرتے ہیں اور قدم دوم مکمل کرتے ہیں۔
 قدم دوم میں ہم فرض کرتے ہیں کہ کسی \عددی{k\ge 7} کے لئے \عددی{k!>3^k} ہے۔ یوں درج ذیل ہو گا۔
\[(k+1)!=(k+1)(k!)>(k+1)3^k>7\cdot 3^k>3^{k+1}\]
لہٰذا  \عددی{k\ge 7} کے لئے  \عددی{k!>3^k} سے مراد \عددی{(k+1)!\ge 3^{k+1}} ہو گا۔ ریاضیاتی ترغیب کا اصول ضمانت دیتا ہے کہ تمام \عددی{n\ge 7} کے لئے \عددی{n!>3^n} ہو گا۔
\انتہا{مثال}

\حصہء{سوالات}

\ابتدا{سوال}
یہ مان لینے کے بعد کہ  ہر دو اعداد \عددی{a} اور \عددی{b}  کے لئے  مثلثی عدم مساوات \عددی{|a+b|\le |a|+|b|}مطمئن  ہو گا، ثابت کریں کہ کسی بھی عدد \عددی{n} کے لئے درج ذیل ہو گا۔
\[|x_1+x_2+\cdots+x_n| \le |x_1|+|x_2|+\cdots +|x_n|\]
\انتہا{سوال}
\ابتدا{سوال}
دکھائیں کہ  \عددی{r\ne 1} کی صورت میں کسی بھی مثبت عدد صحیح \عددی{n} کے لئے درج ذیل ہو گا۔
\[1+r+r^2+\cdots+r^n=\frac{1-r^{n+1}}{1-r}\]
\انتہا{سوال}
\ابتدا{سوال}
حاصل ضرب کا قاعدہ \عددی{\frac{\dif}{\dif x}(uv)=u\frac{\dif v}{\dif x}+v\frac{\dif u}{\dif x}} اور  یہ جانتے ہوئے کہ \عددی{\frac{\dif}{\dif x}(x)=1} ہے دکھائیں کہ ہر مثبت عدد صحیح \عددی{n} کے لئے \عددی{\frac{\dif}{\dif x}(x^n)=nx^{n-1}} ہو گا۔
\انتہا{سوال}
\ابتدا{سوال}
فرض کریں تفاعل \عددی{f(x)} کی خاصیت یہ ہے کہ  کسی بھی دو عدد صحیح \عددی{x_1} اور \عددی{x_2} کے لئے \عددی{{f(x_1x_2)=f(x_1)+f(x_2)}}  ہوتا ہے۔ دکھائیں کہ  کسی بھی \عددی{n} مثبت اعداد \عددی{x_1}، \عددی{x_2}،\نقطے، \عددی{x_n} کے لئے \عددی{f(x_1x_2\cdots x_n)=f(x_1)+f(x_2)+\cdots+f(x_n)} ہو گا۔
\انتہا{سوال}
\ابتدا{سوال}
دکھائیں کہ تمام مثبت   اعداد \عددی{n} کے لئے درج ذیل ہو گا۔
\[\frac{2}{3^1}+\frac{2}{3^2}+\cdots+\frac{2}{3^n}=1-\frac{1}{3^n}\]
\انتہا{سوال}
\ابتدا{سوال}
دکھائیں کہ تمام بڑے  اعداد \عددی{n} کے لئے \عددی{2^n>n^2} ہو گا۔
\انتہا{سوال}
\ابتدا{سوال}
دکھائیں کہ تمام بڑے  اعداد \عددی{n} کے لئے \عددی{2^n\ge n^2} ہو گا۔
\انتہا{سوال}
\ابتدا{سوال}
دکھائیں کہ \عددی{n\ge -3} کے لئے \عددی{2^n\ge 1\!/\!8} ہو گا۔
\انتہا{سوال}
\ابتدا{سوال}\موٹا{مربعوں کے مجموعے}\quad 
دکھائیں کہ پہلے \عددی{n} مثبت عدد صحیح کے مربعوں کا مجموعہ درج ذیل ہو گا۔
\[\frac{n(n+\frac{1}{2})(n+1)}{3}\]
\انتہا{سوال}
\ابتدا{سوال}\موٹا{مکعب کے مجموعے}\quad 
دکھائیں کہ پہلے \عددی{n} مثبت عدد صحیح کے مربعوں کا مجموعہ \عددی{(n(n+1)\!/\!2)^2}  ہو گا۔
\انتہا{سوال}
\ابتدا{سوال}\موٹا{محدود مجموعوں کے قواعد}\quad
دکھائیں کہ محدود مجموعوں کے  درج ذیل  قواعد ہر مثبت عدد صحیح \عددی{n} کے لئے مطمئن ہوں گے۔
\begin{enumerate}[a.]
\item
\[\sum_{k=1}^{n}(a_k+b_k)=\sum_{k=1}^{n}a_k+\sum_{k=1}^{n}b_k\]
\item
\[\sum_{k=1}^{n}(a_k-b_k)=\sum_{k=1}^{n}a_k-\sum_{k=1}^{n}b_k\]
\item
\[\sum_{k=1}^{n} ca_k=c\cdot \sum_{k=1}^{n} a_k \quad \text{\RL{($c$ کوئی بھی عدد ہو سکتا ہے)}}\]
\item
\[\sum_{k=1}^{n} a_k=n\cdot c\quad \text{\RL{(اگر $a_k$ کی قیمت مستقل $c$ ہو)}}\]
\end{enumerate}
\انتہا{سوال}
\ابتدا{سوال}
دکھائیں کہ ہر مثبت عدد صحیح \عددی{n} اور ہر حقیقی عدد \عددی{x} کے لئے \عددی{|x^n|=|x|^n} ہو گا۔
\انتہا{سوال}

